\documentclass[12pt,letterpaper]{extarticle}
\newcommand{\packetlevel}{Grades 4–5}
\newcommand{\packetid}{F03-U-v5}
% Shared layout for the Week 3 packets.  Each packet defines \packetlevel and \packetid first.
\usepackage[T1]{fontenc}
\usepackage[utf8]{inputenc}
\usepackage[default]{sourcesanspro}
\usepackage[letterpaper, left=0.5in, right=0.5in, top=0.72in, bottom=0.68in,
            headheight=18pt, headsep=0.2in, footskip=0.34in]{geometry}
\usepackage{xcolor}
\usepackage{tikz}
\usetikzlibrary{arrows.meta}
\usepackage{fancyhdr}
\usepackage{array}

\definecolor{pbgreen}{HTML}{3FAE4A}
\definecolor{pbblue}{HTML}{2F6FD6}
\definecolor{pbred}{HTML}{D8343C}
\definecolor{pbyellow}{HTML}{F5C518}
\definecolor{pbpurple}{HTML}{7D46A8}
\definecolor{pbpink}{HTML}{F07AB4}
\definecolor{pbteal}{HTML}{1EA59B}
\definecolor{pbgray}{HTML}{8C9099}
\definecolor{slotfill}{gray}{0.975}
\definecolor{framecol}{gray}{0.6}

\pagestyle{fancy}
\fancyhf{}
\fancyhead[L]{\normalsize Week 3 / Shuffle machines / \packetlevel}
\fancyfoot[L]{\small Bellingham Math Circle / Week 3 / \packetid}
\fancyfoot[R]{\small \thepage}
\renewcommand{\headrulewidth}{0pt}
\renewcommand{\footrulewidth}{0pt}

\setlength{\parindent}{0pt}
\setlength{\parskip}{0pt}
\raggedright
\pagenumbering{arabic}

\newcommand{\labelfont}{\small}
\newcommand{\cardfont}{\normalsize\bfseries}

\newcounter{prob}
\newcommand{\problem}[1]{\stepcounter{prob}\par\textbf{Problem \theprob:} #1\par}
% a diagram, centred, with space above it
\newcommand{\fig}[2][10pt]{\par\vspace{#1}{\centering\input{fig/#2}\par}}


\begin{document}

One turn: move every block down its own arrow, then slide the whole bottom row straight up into the top row.
When you draw a machine, give every top slot one arrow going out and every bottom slot one arrow coming in.
A block is home when it is back under its own picture.

\vspace{8pt}
\problem{Put each block under its picture and run each machine. For each block, find the first turn after which it is back home. Then find the first turn after which all the blocks are back home at the same time.}

\fig{u-p1a}
\fig[12pt]{u-p1b}

\newpage
\problem{Without moving any blocks, find how many turns each of these machines takes to bring all the blocks home.}

\fig[16pt]{u-p2}

\newpage
\problem{Which numbers of turns can a machine with 5 slots take to bring all the blocks home? Draw a machine for each number you find, and explain why no other number is possible.}

\fig[16pt]{u-p3}

\newpage
\problem{Find a machine with 6 slots that takes as many turns as possible to bring all the blocks home. Do the same for 7 slots. For each, explain why no machine with that many slots takes more turns.}

\fig[16pt]{u-p4}

\newpage
\problem{A perfect shuffle splits a deck into a top half and a bottom half of the same size, then puts the halves together one card from each half at a time, keeping the top card on top. Do the shuffles with the cards laid out in a row, top card on the left. With the eight cards A, 2, 3, 4, 5, 6, 7, 8 in order from the top, one perfect shuffle gives A, 5, 2, 6, 3, 7, 4, 8. How many perfect shuffles bring these eight cards back to their starting order? Find the number for 4, 6, 10 and 12 cards too, using the cards from A upward in one suit. Then draw the perfect shuffle of 8 cards as a machine.}

\vspace{18pt}
{\centering\renewcommand{\arraystretch}{2.2}%
\begin{tabular}{|l|*{5}{>{\centering\arraybackslash}p{0.75in}|}}
\hline
cards & 4 & 6 & 8 & 10 & 12 \\ \hline
shuffles & & & & & \\ \hline
\end{tabular}\par}

\fig[36pt]{u-p5}

\newpage
\problem{One turn of machine A followed by one turn of machine B moves the blocks the same way as one turn of a single machine, called A then B. For each pair of machines A and B, draw A then B and B then A.}

\fig[16pt]{u-p6}

\newpage
\problem{A machine undoes A if A then that machine puts every block home. For each pair of machines A and B below, draw A then B, and draw the machines that undo A, B and A then B. Explain how to get the machine that undoes A then B from the machines that undo A and B, for any two machines A and B.}

\fig[12pt]{u-p7}

\newpage
\problem{Find two different machines A and B with 4 slots, each moving all four blocks, so that A then B is the same machine as B then A.}

\fig[14pt]{u-p8}

\vspace{30pt}
\problem{For each number of slots from 1 to 10, find the largest number of turns a machine can take to bring all the blocks home. Can you always build a slower machine when you have one more slot? Explain.}

\vspace{18pt}
{\centering\renewcommand{\arraystretch}{2.2}%
\begin{tabular}{|l|*{10}{>{\centering\arraybackslash}p{0.42in}|}}
\hline
slots & 1 & 2 & 3 & 4 & 5 & 6 & 7 & 8 & 9 & 10 \\ \hline
most turns & & & & & & & & & & \\ \hline
\end{tabular}\par}

\vspace{2.4in}
\problem{How many perfect shuffles bring a full deck of 52 cards back to its starting order? Explain how you know.}

\end{document}
